\section{Which reachability claims survive on infinite graphs?}\label{sec:infinite}
A first-order theory of numbers can speak about the graph of a specified
PR relation on natural labels. A relation program in the next result is fixed
in advance. Its numerical parameters may vary. This covers some countably
infinite graphs while keeping the interpretation first-order.

\begin{proposition}[Paths of specified length in a bounded neighbourhood]\leavevmode\label{idx:infinite-paths-of-specified-length-in-a-bounded-neighbourhood}
Suppose adjacency $E(x,y)$ is PR and every successor $y$ of $x$ satisfies
$y\le b(x)$ for a specified PR function $b$. Given $s,t,k$ there is a PR
decision, with witness, for the existence of a walk from $s$ to $t$ of length
at most $k$. A given finite walk can be shortened to a simple path of no
greater length. An odd closed walk in an undirected graph contains an odd cycle.
\end{proposition}
\begin{proof}
Set $M_0=s$ and $M_{j+1}=\max(M_j,\max_{x\le M_j}b(x))$. By induction on
position, every walk of length at most $k$ starting at $s$ stays in $[M_k+1]$.
Enumerate all lengths up to $k$ and the corresponding vertex words in
radix $\max(2,M_k+1)$. Take the first accepted word. Erasure of the segment
between two equal vertices proves shortening exactly as in
Theorem~\ref{stshort}, using the length of the input walk for fuel.
For the last assertion, whenever a closed walk of odd length has a repeated
vertex other than its required final return, split there into two shorter
closed walks. Exactly one length is odd. Continue with that one for at most
the initial length. The final odd closed walk with no intermediate repetition
is an odd cycle (a length-one loop is allowed).
\end{proof}

\begin{proposition}[A uniform obstruction to deciding arbitrary reachability]\leavevmode\label{idx:infinite-a-uniform-obstruction-to-deciding-arbitrary-reachabilit}
There is a single PR directed acyclic graph on natural-number labels, of
undirected degree at most three, for which deciding reachability between
arbitrary designated pairs would decide the halting problem. The equivalence
for each input is provable in iPRA.
\end{proposition}
\begin{proof}
Use a fixed enumeration of deterministic machines with a halting state.
Let $H(e,j)$ say that machine $e$ on input $e$ has halted by step $j$.
This is PR: read the finite instruction table by a bounded scan, store the
finite tape on the two sides of the head as two base-$b$ stacks, and implement
one transition by quotient, remainder and appending a digit. Encode the state,
two stacks and head symbol in a tuple, make the halt state absorbing, and
iterate $j$ times. Invalid descriptions can stand for a fixed looping machine.
Write $H_0(e,j)$ for $H(e,j)$ together with $j=0$ or $H(e,j-1)=0$.
One layer per machine contains a ray $v_{e,0}\to v_{e,1}\to\cdots$ and a
distinct sink $z_e$. Add $v_{e,j}\to z_e$ exactly when $H_0(e,j)$.
Pairing with tags gives PR adjacency and separates layers. There is at most
one added edge in a layer. Hence the degree bound holds and cycles are
impossible. Existence of a walk $v_{e,0}\leadsto z_e$ is equivalent to
$\exists j\,H(e,j)$. From $j$ find the first successful step at most $j$ and
write its finite ray segment plus the final edge. From a walk to the sink read
the time on its last predecessor. The path test and both conversions are PR
in the supplied witnesses. These manipulations of existential witnesses prove
the two implications intuitionistically.

For the metamathematical consequence, a decision procedure for the designated
reachability test would give a machine $D$ which on input $e$ loops when the
decision says reachable, and halts otherwise. Apply the purported decision to
the index of $D$ on itself. Either returned bit contradicts its defining
behaviour. Thus the finite-graph decision theorem cannot be extended to all
such graphs by merely supplying a PR adjacency predicate and a degree bound.
\end{proof}
